% capitulos/01-introducao.tex

Lorem ipsum dolor sit amet, consectetur adipiscing elit. Pellentesque
habitant morbi tristique senectus et netus et malesuada fames ac turpis
egestas \cite{exemplo:artigo:2024}. O estudo de ecossistemas do Cerrado
brasileiro tem ganhado crescente atenção científica em razão da
singularidade de sua biodiversidade e da acelerada degradação a que está
sujeito \cite{exemplo:livro:2020}.

A contaminação ambiental por metais pesados representa um problema de
saúde pública e ambiental de escala global. Conforme
\citeonline{exemplo:tese:2023}, as concentrações de mercúrio em solos
adjacentes a áreas de garimpo ilegal têm excedido consistentemente os
valores de referência estabelecidos pelos órgãos reguladores.

\begin{figure}[H]
  \centering
  \fbox{\rule{0pt}{5cm}\rule{10cm}{0pt}}
  \caption{Mapa de localização da área de estudo no bioma Cerrado,
           Distrito Federal, Brasil. Fonte: elaborado pelo autor (2025).}
  \label{fig:area-estudo}
\end{figure}

A \autoref{fig:area-estudo} ilustra a área de estudo selecionada.
A escolha da região foi motivada pela presença de formações geológicas
com histórico de atividade mineradora e pela proximidade a unidades
de conservação federais \cite{exemplo:web:2024}.

\section{Contexto e Justificativa}

Ut enim ad minim veniam, quis nostrud exercitation ullamco laboris nisi
ut aliquip ex ea commodo consequat \cite{exemplo:artigo:2024}. Duis aute
irure dolor in reprehenderit in voluptate velit esse cillum dolore eu
fugiat nulla pariatur. Excepteur sint occaecat cupidatat non proident,
sunt in culpa qui officia deserunt mollit anim id est laborum.

A relevância científica desta pesquisa reside na lacuna identificada na
literatura quanto aos mecanismos de transporte e acumulação de
contaminantes em gradientes topográficos de savana
\cite{exemplo:capitulo:2019}. Estudos anteriores \cite{exemplo:anais:2023}
apontam que a dinâmica hidrológica do Cerrado cria condicionantes
biogeoquímicos ainda pouco compreendidos.

\section{Hipótese}

A hipótese central deste trabalho é que as concentrações de mercúrio
total nos solos e na serapilheira variam significativamente em função
da proximidade a vias de tráfego intenso e à profundidade do lençol
freático, sendo moduladas pelo teor de matéria orgânica do solo
(\abreviatura{MOS}{Matéria Orgânica do Solo}).

\section{Estrutura do Trabalho}

Este trabalho está organizado em seis capítulos. O
\autoref{cap:metodologia} descreve os materiais e métodos empregados.
O \autoref{cap:recursos} apresenta o orçamento e as fontes de
financiamento. O \autoref{cap:cronograma} detalha o cronograma de
execução. O \autoref{cap:resultados} apresenta e discute os resultados.
Por fim, o \autoref{cap:conclusao} traz as considerações finais.
